\chapter{Fundamentação Teórica}
\label{cap:fundamentacao}

\section{Criptografia Baseada em Identidade}
\label{sec:conceitos}

\subsection{Definição Formal}

Um esquema IBE é formado por quatro algoritmos~\cite{boneh2001ibe}:

\begin{itemize}
  \item \textbf{Setup}$(1^\lambda)$: recebe o parâmetro de segurança $\lambda$ e produz os
    parâmetros públicos \textit{params} e o segredo mestre $\mathit{msk}$. Os parâmetros são
    publicados; o segredo fica com o PKG.

  \item \textbf{Extract}$(\mathit{params}, \mathit{msk}, \mathit{ID})$: dado um identificador
    $\mathit{ID} \in \{0,1\}^*$, o PKG computa e entrega ao usuário sua chave privada
    $d_{\mathit{ID}}$, após verificar sua identidade.

  \item \textbf{Encrypt}$(\mathit{params}, \mathit{ID}, M)$: cifra a mensagem $M$ para o
    destinatário com identidade $\mathit{ID}$, produzindo o texto cifrado $C$.

  \item \textbf{Decrypt}$(\mathit{params}, d_{\mathit{ID}}, C)$: usando a chave privada
    $d_{\mathit{ID}}$, decifra $C$ e recupera $M$, ou retorna $\perp$ em caso de falha.
\end{itemize}

O ponto central é que qualquer um pode cifrar para uma identidade $\mathit{ID}$ conhecendo
apenas os parâmetros públicos, sem certificados e sem contato prévio com o
destinatário~\cite{franklin2008nist}.

\subsection{Noções de Segurança}

A segurança do IBE é avaliada em dois níveis~\cite{boneh2001ibe}:

\begin{itemize}
  \item \textbf{IND-ID-CPA}: o adversário pode consultar chaves privadas de quaisquer identidades,
    exceto a identidade-alvo $\mathit{ID}^*$, e recebe o texto cifrado de uma das duas mensagens de
    sua escolha. Sua vantagem em distingui-las deve ser desprezível. O esquema BasicIdent de Boneh
    e Franklin satisfaz essa noção.

  \item \textbf{IND-ID-CCA}: além das consultas de chaves, o adversário pode submeter textos cifrados a
    um oráculo de decifração (exceto o texto desafio). Essa noção mais forte é satisfeita pelo
    FullIdent, por meio de um teste de rejeição que impede a aceitação de textos cifrados modificados.
\end{itemize}

A diferença entre BasicIdent e FullIdent está na adição de duas funções de hash auxiliares
($H_3$ e $H_4$) e de um teste de consistência na decifração, a \textit{Fujisaki-Okamoto
transform}~\cite{fujisaki1999secure,boneh2001ibe}.

O esquema BasicIdent apenas faz XOR da mensagem com $H_2(\hat{e}(Q_{\mathit{ID}},
P_{\mathit{pub}})^r)$, o que torna o texto cifrado maleável. Um atacante que intercepta
$C = \langle U, V \rangle$ pode inverter um bit qualquer de $V$ e obter um novo texto cifrado
válido, cuja decifração é a mensagem original com o bit correspondente invertido. Se o
atacante tiver acesso a um oráculo de decifração, como permite o jogo IND-ID-CCA, ele envia
esse texto cifrado modificado, recebe de volta a mensagem alterada e desfaz a inversão para
recuperar a mensagem original, sem nunca ter obtido a chave privada~\cite{menezes2009intro}.
Esse é precisamente o ataque que o FullIdent bloqueia: qualquer alteração em $U$, $V$ ou $W$
faz com que o $\sigma'$ recuperado na decifração seja diferente do $\sigma$ original, o que
propaga um $r'$ inconsistente e faz a verificação $U = r' \cdot P$ falhar (ver
Seção~\ref{sec:fullident}). O oráculo de decifração deixa de ser útil ao atacante porque ele só
responde a textos cifrados que já eram válidos.

\subsection{Comparação com PKI e Custódia de Chaves}

Em relação à PKI, o IBE não precisa de Autoridades Certificadoras nem
de certificados. Qualquer um pode cifrar para uma identidade mesmo que o destinatário ainda não
tenha se registrado, já que as mensagens ficam esperando até que ele obtenha sua chave
privada junto ao PKG~\cite{franklin2008nist}.

O lado negativo é que o PKG conhece o segredo mestre $s$ e pode calcular a chave privada de
qualquer usuário e, com isso, ler qualquer mensagem do sistema. Esse problema se chama
\textit{custódia de chaves} (\textit{key escrow}) e é uma limitação do modelo IBE em
geral~\cite{franklin2008nist}.

% ---

\section{Fundamentos Matemáticos}
\label{sec:funcionamento}

\subsection{Curvas Elípticas e Grupos}

Uma curva elíptica $E$ sobre um corpo finito $\mathbb{F}_p$ ($p$ primo) é o conjunto de pontos
$(x, y) \in \mathbb{F}_p^2$ satisfazendo uma equação da forma $y^2 = x^3 + ax + b$, juntamente
com o ponto no infinito $\mathcal{O}$~\cite{menezes2009intro}. Os pontos de $E(\mathbb{F}_p)$
formam um grupo abeliano com a lei de adição de pontos. O esquema FullIdent utiliza uma curva
supersingular do \textbf{Tipo A}: $y^2 = x^3 + x$ sobre $\mathbb{F}_p$, com grau de imersão
$k = 2$, o que significa que o grupo de chegada do emparelhamento vive em $\mathbb{F}_{p^2}$.

O grupo $G_1$ é um subgrupo cíclico de $E(\mathbb{F}_p)$ de ordem prima $q$, gerado por um
ponto $P$. O grupo alvo é $G_T$, um subgrupo multiplicativo de ordem $q$ em
$\mathbb{F}_{p^2}^*$.

\subsection{Emparelhamentos Bilineares}

Um \textbf{emparelhamento bilinear} é uma aplicação $\hat{e}: G_1 \times G_1 \rightarrow G_T$
que satisfaz as seguintes propriedades~\cite{menezes2009intro,boneh2001ibe}:

\begin{itemize}
  \item \textbf{Bilinearidade}: para todos $P, Q \in G_1$ e $a, b \in \mathbb{Z}$:
    \[ \hat{e}(aP,\, bQ) = \hat{e}(P,\, Q)^{ab}. \]

  \item \textbf{Não-degenerescência}: existe $P \in G_1$ tal que $\hat{e}(P, P) \neq 1_{G_T}$.

  \item \textbf{Computabilidade}: existe um algoritmo eficiente para calcular $\hat{e}(P, Q)$
    para quaisquer $P, Q \in G_1$.
\end{itemize}

A construção original de Boneh-Franklin usa o emparelhamento de Weil modificado sobre curvas
supersingulares~\cite{boneh2001ibe}. Na prática, porém, o emparelhamento de Tate é preferido:
os próprios Boneh e Franklin observam que tanto a cifragem quanto a decifragem do FullIdent
podem ser aceleradas usando Tate em vez de Weil~\cite{boneh2001ibe}, e é essa família de
emparelhamentos (Tate e suas variantes mais rápidas, como o emparelhamento $\eta$) que a
literatura de implementação recomenda e que bibliotecas como a PBC de fato calculam para
curvas do Tipo A~\cite{menezes2009intro}. O grau de imersão baixo ($k=2$) torna esse cálculo
eficiente, sem comprometer a dificuldade do logaritmo discreto em $G_1$~\cite{menezes2009intro}.

Neste trabalho, usamos a biblioteca PBC (\textit{Pairing-Based Cryptography})~\cite{pbc} para
os emparelhamentos, com os parâmetros padrão \texttt{a.param} da própria biblioteca.

\subsection{O Problema Bilinear de Diffie-Hellman (BDH)}

O FullIdent não se apoia diretamente no problema computacional de Diffie-Hellman (CDH) em
$G_1$, como fazem os esquemas clássicos baseados em curvas elípticas, porque o próprio
emparelhamento bilinear destrói a dificuldade do problema \emph{decisional} de Diffie-Hellman
(DDH) nesse grupo. Dado $(P, aP, bP, cP)$, basta testar se $\hat{e}(P, cP) = \hat{e}(aP, bP)$
para decidir se $c \equiv ab \pmod q$, já que por bilinearidade $\hat{e}(P, cP) =
\hat{e}(P,P)^c$ e $\hat{e}(aP, bP) = \hat{e}(P,P)^{ab}$. Assim, o DDH em $G_1$ passa a ser
fácil de resolver assim que existe um emparelhamento eficiente sobre esse
grupo~\cite{boneh2001ibe,menezes2009intro}, e esquemas como o Cramer-Shoup, que dependem da
dificuldade de DDH para atingir segurança CCA, não podem ser usados diretamente sobre $G_1$
por esse motivo. O CDH em $G_1$ (encontrar $abP$) segue acreditado como difícil, mas isso não é
suficiente para as reduções de segurança do FullIdent, que passam pelo emparelhamento;
por isso Boneh e Franklin definem um problema novo, o BDH, adaptado exatamente às
propriedades do emparelhamento.

O problema BDH é definido como~\cite{boneh2001ibe}: dado $(P, aP, bP, cP) \in G_1^4$, com
$a, b, c \in \mathbb{Z}_q^*$ escolhidos uniformemente ao acaso, calcular $\hat{e}(P, P)^{abc}
\in G_T$.

A hipótese BDH afirma que nenhum algoritmo probabilístico de tempo polinomial consegue resolver
esse problema com vantagem não desprezível. A segurança do FullIdent depende diretamente dessa
hipótese. Qualquer adversário que quebre o esquema pode ser convertido em um algoritmo que
resolve BDH~\cite{boneh2001ibe}.

% ---

\section{O Esquema FullIdent}
\label{sec:fullident}

\subsection{Funções de Hash}

FullIdent utiliza quatro funções de hash~\cite{boneh2001ibe}, sendo $n$ o comprimento em bits
da mensagem:

\begin{align*}
  H_1 &: \{0,1\}^* \rightarrow G_1^* \\
  H_2 &: G_T \rightarrow \{0,1\}^n \\
  H_3 &: \{0,1\}^n \times \{0,1\}^n \rightarrow \mathbb{Z}_q^* \\
  H_4 &: \{0,1\}^n \rightarrow \{0,1\}^n
\end{align*}

Neste trabalho, para $H_1$ usamos \texttt{element\_from\_hash} da biblioteca PBC (que
internamente usa SHA-256 para mapear a identidade a um ponto de curva); $H_2$, $H_3$ e $H_4$
são implementadas sobre SHA-256 com separação de domínio.

\subsection{Algoritmos}

\paragraph{Setup.} O PKG escolhe um grupo $G_1$ de ordem $q$ com gerador $P$, um escalar
aleatório $s \in \mathbb{Z}_q^*$ (segredo mestre) e computa:
\[ P_{\mathit{pub}} = s \cdot P \quad \in G_1. \]
Os parâmetros públicos são $\mathit{params} = (G_1, G_T, \hat{e}, n, P, P_{\mathit{pub}},
H_1, H_2, H_3, H_4)$.

\paragraph{Extract.} Para uma identidade $\mathit{ID} \in \{0,1\}^*$, após autenticação do
usuário, o PKG computa:
\[ Q_{\mathit{ID}} = H_1(\mathit{ID}) \in G_1^*, \qquad d_{\mathit{ID}} = s \cdot Q_{\mathit{ID}} \in G_1^*. \]

\paragraph{Encrypt.} Para cifrar $M \in \{0,1\}^n$ para o destinatário $\mathit{ID}$:
\begin{enumerate}
  \item Compute $Q_{\mathit{ID}} = H_1(\mathit{ID})$.
  \item Escolha $\sigma \in \{0,1\}^n$ uniformemente ao acaso.
  \item Compute $r = H_3(\sigma, M)$.
  \item Compute as três componentes do texto cifrado $C = \langle U, V, W \rangle$:
    \begin{align}
      U &= r \cdot P, \label{eq:U}\\
      V &= \sigma \oplus H_2\!\left(\hat{e}(Q_{\mathit{ID}},\, P_{\mathit{pub}})^r\right), \label{eq:V}\\
      W &= M \oplus H_4(\sigma). \label{eq:W}
    \end{align}
\end{enumerate}

\paragraph{Decrypt.} Para decifrar $C = \langle U, V, W \rangle$ com chave privada
$d_{\mathit{ID}}$:
\begin{enumerate}
  \item Compute $\sigma' = V \oplus H_2(\hat{e}(d_{\mathit{ID}},\, U))$.
  \item Compute $M' = W \oplus H_4(\sigma')$.
  \item Compute $r' = H_3(\sigma', M')$.
  \item \textbf{Verificação de consistência:} se $U \neq r' \cdot P$, retorne $\perp$.
  \item Caso contrário, retorne $M'$.
\end{enumerate}

\subsection{Corretude}

A decifração recupera $M$ corretamente por causa da bilinearidade do emparelhamento. Para um
texto cifrado gerado com escalar $r$ e para o destinatário com chave privada
$d_{\mathit{ID}} = s \cdot Q_{\mathit{ID}}$:

\begin{align*}
  \hat{e}(d_{\mathit{ID}},\, U)
  &= \hat{e}(s \cdot Q_{\mathit{ID}},\, r \cdot P) \\
  &= \hat{e}(Q_{\mathit{ID}},\, P)^{sr} \\
  &= \hat{e}(Q_{\mathit{ID}},\, s \cdot P)^r \\
  &= \hat{e}(Q_{\mathit{ID}},\, P_{\mathit{pub}})^r.
\end{align*}

Portanto, $H_2(\hat{e}(d_{\mathit{ID}}, U)) = H_2(\hat{e}(Q_{\mathit{ID}},
P_{\mathit{pub}})^r)$, e a XOR com $V$ recupera $\sigma$ corretamente; a XOR de $W$ com
$H_4(\sigma)$ recupera $M$; e a verificação $U = r' \cdot P$ é satisfeita.

O teste de rejeição no passo 4 é o que impede ataques de texto cifrado escolhido. Qualquer
modificação em $U$, $V$ ou $W$ produz um $\sigma'$ inválido, que por sua vez gera $r'$ e $M'$
inconsistentes, fazendo com que $r' \cdot P \neq U$, e a decifração é
abortada~\cite{boneh2001ibe}.

\subsection{Identidades Temporais e Revogação}

O IBE tem suporte natural a \textbf{identidades temporais}~\cite{franklin2008nist}. Em vez de
usar \texttt{alice@ufsc.br} como identidade, é possível usar \texttt{alice@ufsc.br|2026-06}.
Qualquer remetente pode cifrar para esse identificador, mas só Alice consegue decifrar se tiver
obtido a chave do período correspondente junto ao PKG. Após o período expirar, a chave torna-se
inútil automaticamente, sem precisar de listas de revogação.

Além das identidades temporais, o PKG pode revogar uma identidade diretamente, passando seu
status para \textit{revoked} no registro interno. Com isso, o PKG bloqueia novas extrações.
Chaves já emitidas continuam válidas até o fim do período, o que é esperado em qualquer
implementação de IBE~\cite{franklin2008nist}.

% ---

\section{Exemplo Ilustrativo}
\label{sec:exemplo}

Para entender melhor como o emparelhamento funciona na prática, vamos a um exemplo com
números pequenos. Os parâmetros aqui são muito menores do que o necessário na prática e servem
só para ilustração.

Seja $G_1$ um subgrupo cíclico de ordem $q = 11$ com gerador $P$. Suponha:
\begin{itemize}
  \item Segredo mestre: $s = 3$, logo $P_{\mathit{pub}} = 3P$.
  \item Identidade do destinatário: $\mathit{ID} = \texttt{bob@ufsc.br}$, com
        $Q_{\mathit{ID}} = H_1(\mathit{ID}) = 5P$ (valor hipotético após o mapeamento).
  \item Chave privada de Bob: $d_{\mathit{ID}} = 3 \cdot 5P = 15P$.
        Como $15 \equiv 4 \pmod{11}$, tem-se $d_{\mathit{ID}} = 4P$.
\end{itemize}

O remetente escolhe $\sigma$ aleatório, computa $r = H_3(\sigma, M)$ e suponha $r = 2$. O
valor de $V$ é obtido com:
\[
  \hat{e}(Q_{\mathit{ID}},\, P_{\mathit{pub}})^r
  = \hat{e}(5P,\, 3P)^2
  = \hat{e}(P,\, P)^{15 \cdot 2}
  = \hat{e}(P,\, P)^{30}
  \equiv \hat{e}(P,\, P)^{8} \pmod{q}.
\]

Na decifração, Bob computa:
\[
  \hat{e}(d_{\mathit{ID}},\, U)
  = \hat{e}(4P,\, 2P)
  = \hat{e}(P,\, P)^{4 \cdot 2}
  = \hat{e}(P,\, P)^{8}.
\]

Os dois valores coincidem ($30 \equiv 8 \pmod{11}$), confirmando que Bob recupera $\sigma$
corretamente. Caso Eve tente decifrar com uma chave $d_{\mathit{Eve}} = 7P$, ela obteria:
\[
  \hat{e}(7P,\, 2P) = \hat{e}(P,\, P)^{14} \equiv \hat{e}(P,\, P)^{3},
\]
que é distinto de $\hat{e}(P,P)^8$. O $\sigma'$ recuperado seria inválido, $r'$ e $M'$
inconsistentes, e a verificação $U \stackrel{?}{=} r' \cdot P$ falharia, retornando $\perp$.
